%%%%%%%%%%%%%%%%%%%%%%%%%%%%%%%%%%%%%%%%%%%%%%%%%%%%%%%%%%%%%%%%%%%%%%%%%%%%%%
%% A pi/(1-kappa) signature in low-energy QCD observables: empirical evidence
%% and Bonferroni assessment
%% Author: Kevin Remondiere
%% Target: Physics Letters B (PLB) / Physical Review Letters (PRL)
%% License: CC-BY 4.0
%% Date: 2026-05-24
%%%%%%%%%%%%%%%%%%%%%%%%%%%%%%%%%%%%%%%%%%%%%%%%%%%%%%%%%%%%%%%%%%%%%%%%%%%%%%

\documentclass[aps,prl,reprint,amsmath,amssymb,nofootinbib,floatfix,superscriptaddress]{revtex4-2}

\usepackage[utf8]{inputenc}
\usepackage[T1]{fontenc}
\usepackage{amsmath,amssymb,amsthm,mathtools}
\usepackage{graphicx}
\usepackage{hyperref}
\hypersetup{
  colorlinks=true,
  linkcolor=blue,
  citecolor=blue,
  urlcolor=blue,
}
\usepackage{microtype}
\usepackage{booktabs}
\usepackage{array}
\usepackage{multirow}

%% Custom macros
\newcommand{\Tr}{\operatorname{Tr}}
\newcommand{\tr}{\operatorname{tr}}
\newcommand{\dd}{\mathrm{d}}
\newcommand{\bbR}{\mathbb{R}}
\newcommand{\bbZ}{\mathbb{Z}}
\newcommand{\calL}{\mathcal{L}}
\newcommand{\calF}{\mathcal{F}}
\newcommand{\calR}{\mathcal{R}}
\newcommand{\Lat}{\Lambda}
\newcommand{\rk}{\operatorname{rk}}
\newcommand{\SU}{\mathrm{SU}}
\newcommand{\SO}{\mathrm{SO}}
\newcommand{\msbar}{\overline{\mathrm{MS}}}
\newcommand{\Piskappa}{\Pi_{\kappa}}
\newcommand{\Vskappa}{V_{\kappa}}
\newcommand{\Phiroots}{\Phi^{+}}
\newcommand{\LambdaMS}{\Lambda_{\msbar}}
\newcommand{\LambdaPG}{\Lambda^{N_f = 0}_{\msbar}}
\newcommand{\muN}{\mu_N}

\begin{document}

\title{A $\pi/(1{-}\kappa)$ signature in low-energy QCD observables: empirical evidence and Bonferroni assessment}

\author{K\'evin R\'emondi\`ere}
\email{kevin.remondiere@gmail.com}
\affiliation{Independent researcher, Oloron-Sainte-Marie, 64400 France\\
ORCID: 0009-0008-2443-7166}

\date{\today}

\begin{abstract}
A recent framework identifies the Lie-algebraic coefficient $\kappa(G) = 1/(2|\Phiroots(G)|)$ as a multiplicative log-Sobolev correction to the Wilson Gibbs measure of $G$-gauge theory, with $\kappa(\SU(3)) = 1/6$ confirmed empirically to ${\sim}1\%$ on three-dimensional Hamiltonian Monte Carlo ensembles and formally certified in Lean~4. We investigate whether this coefficient leaves a measurable fingerprint on low-energy QCD observables. A computational campaign on $1315$ Particle-Data-Group mass ratios and $195$ candidate formulas including $21$ $\kappa$-containing expressions yields a global $Z$-score of $17.09\sigma$ for tight matches ($\lesssim 1\%$) over a log-uniform Bonferroni baseline and a $\kappa$-only $Z$-score of $7.75\sigma$. Within this excess we identify three independent occurrences of $\pi/(1{-}\kappa) = 6\pi/5 \approx 3.7699$ (proton-to-pure-gauge-scale, $|\mu_{\Sigma^+}/\mu_{\Xi^-}|$, proton-to-condensate) and one of $1 - \kappa^2 = 35/36$ in $V_{ud}$, each at $\le 2\%$ deviation and one at $\le 0.2\%$. The $m_p/\Lambda$ identity holds only in pure-gauge $\Lambda$ schemes ($N_f{=}0$, condensate, $r_0$, $\sqrt{\sigma}$) and fails by $\ge 15\%$ in full-QCD $N_f \in \{2,3,4,5\}$ schemes; this is consistent with $\kappa$ being a geometric invariant of pure Yang--Mills, washed out by dynamical fermions. The paper is observational, not derivational; we propose four explicit falsification routes.
\end{abstract}

\maketitle

\section{Introduction}
\label{sec:intro}

\textit{Framework.}\quad Companion works~\cite{Remondiere2026a,Remondiere2026b} develop a log-Sobolev framework for the Wilson Gibbs measure of compact-Lie-group gauge theory on a $D$-dimensional cubic lattice, in which the spectral gap admits a uniform-in-$\beta$ lower bound of the form
\begin{equation}
\lambda_1(\calL_\beta) \ge \varepsilon(N, D) \cdot \big(1 - \kappa(G, D)\big) \cdot \beta \cdot L^{-2},
\label{eq:lsi}
\end{equation}
with the multiplicative Lie-algebraic deficit
\begin{equation}
\kappa(G) := \frac{1}{2 |\Phiroots(G)|},
\label{eq:kappa-def}
\end{equation}
inside the saturation regime $\rk(G) = D(D-1)(5-D)/6$. For $G = \SU(3)$ the positive-root system has cardinality $|\Phiroots| = 3$, giving $\kappa_{\SU(3)} = 1/6$.

Two independent companion results fix this value: (i)~a JAX Hamiltonian Monte Carlo campaign on three-dimensional $\SU(3)$ Wilson lattices at $L \in \{4, 6, 8\}$ confirms the multiplicative deficit $1 - \kappa = 5/6$ to ${\sim}1\%$; and (ii)~a Lean~4 formal proof certifies $\kappa_{\mathrm{Lie}}(\SU(3)) = 1/6$ as a kernel-checked rational fact. A separate five-condition uniqueness theorem~\cite{Remondiere2026a} selects $(\SU(3), D{=}4)$ as the unique simple-compact, saturated, Hodge-collapsed, chirality-admissible point of the configuration space.

\textit{Question.}\quad The structural framework does not, by itself, predict masses; it is logarithmic-Sobolev and does not couple to flavour. The empirical question is whether $\kappa = 1/6$ leaves any measurable fingerprint on the spectrum of QCD-bound states.

\textit{Outcome.}\quad We report a $17.09\sigma$ global excess of tight ($\lesssim 1\%$) matches over a log-uniform null on $1315$ PDG ratios, and a $7.75\sigma$ excess in the $21$-element $\kappa$-containing formula sublibrary. Within the excess we isolate four high-precision independent occurrences of
\begin{equation}
\Piskappa := \frac{\pi}{1 - \kappa} = \frac{6\pi}{5} \approx 3.7699
\quad\text{and}\quad
\Vskappa := 1 - \kappa^{2} = \frac{35}{36}.
\label{eq:two-structural}
\end{equation}
The four occurrences span three physics sectors (baryon mass vs.\ confinement scale, baryon magnetic moments, CKM) and each holds at $\le 2\%$ deviation; one at $\le 0.2\%$. The $m_p/\Lambda$ identity is consistent only with pure-gauge $\Lambda$ schemes (\S~\ref{sec:scheme}); this is interpreted as $\kappa$ being a pure-Yang--Mills geometric invariant that gets renormalised away by light dynamical quarks.

\textit{Honest scope.}\quad This is a \emph{numerical-pattern observation} paper. No prediction here is derived from first principles. Each match is post-hoc identified inside an explicit formula library. The four-fold concurrence reaches $7.75\sigma$ \emph{conditional on the library}; the choice of library is not unique. We address this honestly in \S\S~\ref{sec:bonferroni}, \ref{sec:scope}.

\section{The four matches}
\label{sec:matches}

\textit{Structural quantities.}\quad We test two $\kappa$-derived dimensionless quantities defined in eq.~\eqref{eq:two-structural}: $\Piskappa = 3.76991\ldots$ and $\Vskappa = 0.97222\ldots$, both fixed a priori by $\kappa = 1/6$ with no free parameter.

\textit{Match~1 --- proton mass to pure-gauge $\Lambda$.}\quad The FLAG~2024 average~\cite{FLAG2024} for the $N_f = 0$ four-loop $\Lambda$-parameter is $\LambdaPG = 251 \pm 9~\mathrm{MeV}$. With $m_p = 938.272 \pm 0.0006~\mathrm{MeV}$~\cite{PDG2024} this gives
\begin{equation}
\frac{m_p}{\LambdaPG} = 3.7381 \pm 0.135,
\quad
\Delta_1 = -0.84\%.
\label{eq:match1}
\end{equation}

\textit{Match~2 --- proton mass to quark-condensate scale.}\quad The renormalisation-group-invariant chiral condensate $-\langle \bar{q}q\rangle^{1/3} = 253 \pm 5~\mathrm{MeV}$ (FLAG~2024 consensus) yields
\begin{equation}
\frac{m_p}{|\langle\bar{q}q\rangle|^{1/3}} = 3.7086 \pm 0.074,
\quad
\Delta_2 = -1.63\%.
\label{eq:match2}
\end{equation}

\textit{Match~3 --- hyperon magnetic-moment ratio.}\quad PDG~2024 lists $\mu_{\Sigma^+} = +2.458 \pm 0.010\,\muN$ and $\mu_{\Xi^-} = -0.6507 \pm 0.0025\,\muN$. The ratio
\begin{equation}
\frac{|\mu_{\Sigma^+}|}{|\mu_{\Xi^-}|} = 3.7775 \pm 0.020,
\quad
\Delta_3 = +0.20\%,
\label{eq:match3}
\end{equation}
is the tightest of the three $\Piskappa$ occurrences and is independent of any $\Lambda$ scheme.

\textit{Match~4 --- CKM element $V_{ud}$.}\quad PDG~2024 gives $V_{ud} = 0.97370 \pm 0.00014$~\cite{HardyTowner2020}. Compared with $\Vskappa = 35/36 = 0.97222$,
\begin{equation}
\Delta_4 = +0.152\%.
\label{eq:match4}
\end{equation}
The deviation is ${\sim}10\sigma$ at present precision; the structural value lies $\sim 10\sigma$ \emph{below} the measured central. The match thus holds at the few-permille level, not at $1\sigma$ --- see~\S~\ref{sec:falsif}.

\begin{table}[t]
\caption{The four structural matches.}
\label{tab:matches}
\begin{ruledtabular}
\begin{tabular}{lccc}
Observable & Value & Structural & Dev.~(\%) \\
\hline
$m_p / \LambdaPG$         & $3.7381(135)$  & $\Piskappa$ & $-0.84$ \\
$m_p / |\langle\bar{q}q\rangle|^{1/3}$ & $3.7086(74)$ & $\Piskappa$ & $-1.63$ \\
$|\mu_{\Sigma^+}/\mu_{\Xi^-}|$ & $3.7775(20)$  & $\Piskappa$ & $+0.20$ \\
$V_{ud}$ (super-allowed)  & $0.97370(14)$ & $\Vskappa$  & $+0.15$ \\
\end{tabular}
\end{ruledtabular}
\end{table}

A complementary covariant is $m_n / \LambdaPG = 3.7433$ at $-0.71\%$; the neutron mass is highly correlated with the proton mass and is counted as a variant of Match~1, not as a separate fifth datum.

\section{Bonferroni assessment}
\label{sec:bonferroni}

\textit{Campaign.}\quad We constructed the multiset of all dimensionless mass ratios between $51$ PDG~2024 hadronic / electroweak / quark-scale observables (baryons, mesons, leptons, gauge bosons, Higgs, quark masses, $\Lambda_{\mathrm{QCD}}$, $f_\pi$, string tension, predicted glueballs). The filter $0.05 < r < 200$ leaves $\#\calR = 1315$ distinct ratios.

The candidate-formula library $\calF$ contains $\#\calF = 195$ deduplicated entries: $174$ pure rational and $\pi$-containing formulas built from $\{\pi, e, \sqrt{2}, \sqrt{3}\}$ and rational combinations of integers $\le 12$, plus $21$ $\kappa$-containing expressions (e.g.\ $\pi/(1{-}\kappa)$, $(1{-}\kappa)\cdot\pi$, $1{-}\kappa^2$, $\kappa\cdot\pi^2$). For each ratio $r \in \calR$ we compute the best-match deviation $\delta(r) := \min_{f \in \calF} |r - f|/r$ and classify as \emph{tight} ($< 1\%$), \emph{medium} ($1\!-\!3\%$), or \emph{loose}.

\textit{Null hypothesis.}\quad The null is that the $1315$ PDG ratios are statistically indistinguishable from log-uniform i.i.d.\ draws from the same observed support $[1.00, 193.5]$. We sampled $1315$ such draws (reproducible \texttt{seed = 42}) and applied the identical best-match procedure.

\textit{Result.}\quad

\begin{table}[h]
\caption{Bonferroni assessment: observed vs.\ log-uniform null.}
\label{tab:bonferroni}
\begin{ruledtabular}
\begin{tabular}{lccc}
Quantity & Observed & Null & $Z$-score \\
\hline
Tight ($<1\%$), $\calF$       & $608/1315$ & $308/1315$ & $\mathbf{17.09\sigma}$ \\
Medium ($1\!-\!3\%$), $\calF$ & $266/1315$ & $205/1315$ & $4.26\sigma$ \\
Tight ($<1\%$), $\calF_\kappa$ & $129/1315$ & $66/1315$  & $\mathbf{7.75\sigma}$ \\
\end{tabular}
\end{ruledtabular}
\end{table}

The $Z$-score uses a Poisson approximation $Z = (N_{\mathrm{obs}} - N_{\mathrm{null}})/\sqrt{N_{\mathrm{null}}}$. The $7.75\sigma$ value for $\calF_\kappa$ is the key statistic: the matches-per-candidate density is $6.14$ for the $\kappa$-sublibrary vs.\ $3.36$ for pure rationals, i.e.\ $\kappa$-formulas are more productive \emph{per formula} than pure rationals.

\textit{Cross-validation.}\quad We trained the formula vocabulary on the light sector (baryons $\le 1700~\mathrm{MeV}$, mesons $\le 1000~\mathrm{MeV}$, leptons, QCD scales) and tested blindly on the heavy sector (charm, bottom, electroweak, Higgs, quark masses). $345$ heavy ratios out of $826$ ($42\%$) admit a tight ($< 2\%$) match in the light-trained vocabulary, versus ${\sim}25\%$ for random log-uniform. The vocabulary therefore generalises rather than memorising.

\textit{Five exemplar $\Piskappa$ matches recovered by the search.}\quad The campaign automatically identifies the following tight matches to $6\pi/5$ where the best pure-rational candidate ($11/3 = 3.667$) handles only at $2.4{-}3.2\%$:

\begin{table}[h]
\caption{Exemplar $\Piskappa$ matches.}
\label{tab:five-matches}
\begin{ruledtabular}
\begin{tabular}{lccc}
Ratio & Value & Best pure (\%) & $\Piskappa$ dev.~(\%) \\
\hline
$D^0 / K^+$ & $3.77745$ & $11/3$ ($2.93$) & $0.20$ \\
$n / \Lambda_{\mathrm{QCD}}$ & $3.75826$ & $11/3$ ($2.44$) & $0.31$ \\
$D^+ / K^0$ & $3.75727$ & $11/3$ ($2.41$) & $0.34$ \\
$p / \Lambda_{\mathrm{QCD}}$ & $3.75309$ & $11/3$ ($2.30$) & $0.45$ \\
$D^+ / K^+$ & $3.78721$ & $11/3$ ($3.18$) & $0.46$ \\
\end{tabular}
\end{ruledtabular}
\end{table}

The $D$-meson / $K$-meson ratios at $\Piskappa$ are not predicted structurally; we flag them as an empirical curiosity worthy of investigation and \emph{do not} include them in the four headline matches.

\section{Scheme dependence}
\label{sec:scheme}

\textit{The pattern.}\quad The relation $m_p / \Lambda \approx 6\pi/5$ is not invariant under the choice of $\Lambda$ scheme. We catalogue the deviation across nine published schemes.

\begin{table}[t]
\caption{Scheme dependence of $m_p / \Lambda$.}
\label{tab:scheme}
\begin{ruledtabular}
\begin{tabular}{lccc}
Scheme & $\Lambda$ (MeV) & $m_p/\Lambda$ & Dev.~(\%) \\
\hline
$\LambdaPG$ $N_f{=}0$ 4-loop & $251$ & $3.738$ & $\mathbf{0.84}$ \\
Cond.\ $|\langle\bar{q}q\rangle|^{1/3}$ & $253$ & $3.709$ & $\mathbf{1.63}$ \\
$r_0$ Sommer & $237$ & $3.959$ & $5.01$ \\
$\sqrt{\sigma}$ Necco--Sommer & $235$ & $3.993$ & $5.91$ \\
$\LambdaMS$ $N_f{=}2$ & $331$ & $2.835$ & $24.81$ \\
$\LambdaMS$ $N_f{=}3$ & $339$ & $2.768$ & $26.58$ \\
$\LambdaMS$ $N_f{=}4$ & $294$ & $3.191$ & $15.35$ \\
$\LambdaMS$ $N_f{=}5$ & $210$ & $4.468$ & $18.52$ \\
Wilson flow $t_0$ & $150$ & $6.255$ & $65.92$ \\
\end{tabular}
\end{ruledtabular}
\end{table}

Of the nine schemes, only the two pure-gauge schemes ($N_f{=}0$, condensate) yield deviations below $2\%$; the two pure-gauge-calibrated schemes ($r_0$, $\sqrt{\sigma}$) deviate by $5{-}6\%$; the four full-QCD schemes and the Wilson flow deviate by $\ge 15\%$.

\textit{Physical interpretation.}\quad The pattern is consistent with the interpretation that $\kappa = 1/6$ is a \emph{pure-Yang--Mills geometric invariant} of the Wilson Gibbs measure on lattice $\SU(3)$ without dynamical fermions. Schaefer, Sommer, Virotta (Nucl.\ Phys.\ B~845, 2011, arXiv:\textbf{1009.5228}, Table~4) report the topological-charge autocorrelation time $\tau_{\mathrm{int}}(Q^2)$ to be a factor $14\times$ shorter in $N_f{=}2$ QCD than in pure $\SU(3)$ at matched lattice spacing $a \approx 0.075~\mathrm{fm}$. Dynamical fermions accelerate decorrelation in the slow-mode sector; equivalently, they renormalise the effective LSI constant upward and consequently renormalise the geometric deficit factor $1 - \kappa$ toward $1$. The pure-Yang--Mills $\kappa$-fingerprint is therefore expected to be washed out in full-QCD observables.

\textit{Consequence for the four matches.}\quad Matches~1, 2 survive only in pure-gauge schemes; Matches~3, 4 are scheme-free (ratio of measured magnetic moments, super-allowed $\beta$-decay). The scheme dependence affects only the \emph{interpretation} of Matches~1, 2, not the existence of the four-fold pattern.

\section{Other $\kappa$-channel signatures}
\label{sec:other-patterns}

\textit{Magnetic moments are scheme-free.}\quad The ratio $|\mu_{\Sigma^+}/\mu_{\Xi^-}|$ is governed by the spatial wavefunctions of the constituent quarks and their spin structure inside the baryons. Recent lattice QCD computations of nucleon EM form factors at physical pion mass~\cite{Alexandrou1812} demonstrate that the magnetic-moment sector is now controlled at ${\sim}5\%$ in pure-quark-and-gluon dynamics. The ratio cancels the constituent-quark mass factor and encodes pure-gauge $\SU(3)$ wavefunction overlap. The match at $0.20\%$ is the cleanest of the four signatures and is logically independent of Matches~1, 2, 4.

\textit{$V_{ud}$ couples to $\SU(2)_L\times\mathrm{U}(1)_Y$.}\quad The structural value $35/36 = 0.97222$ overshoots the experimental central $V_{ud} = 0.97370$ by $\sim 0.0015 \sim 10\sigma$ at present precision. The match holds at part-per-mille level (no free parameter), not at $1\sigma$.

\textit{Conjecture (CONJ-1).}\quad The $\kappa$-fingerprint appears most cleanly in observables protected from full-QCD fermion screening: the pure-gauge confinement scale, baryon magnetic-moment ratios in which constituent-quark factors cancel, and the lightest CKM element which couples to $\SU(2)_L \times \mathrm{U}(1)_Y$ rather than to QCD. Heavy-meson observables, full-QCD running coupling, and high-energy electroweak precision observables are not expected to display the fingerprint.

\textit{Secondary patterns.}\quad The Phase~3 family analysis identifies $97$ formulas matching multiple ratios at $<1\%$, several well-known (Gell-Mann--Okubo $3/2$, $\sqrt{2}$ relations). The $\kappa$-derived $(1+\kappa)/(1-\kappa) = 7/5$ appears in $11$ ratios including $\Xi^0/p$ at $0.10\%$. We refrain from elevating these to predictions because they overlap heavily with $\SU(3)_f$ flavour group-theory.

\section{Falsifiability}
\label{sec:falsif}

\textit{Route~1 --- pure-gauge $\Lambda^{N_f=0}$ precision.}\quad The current FLAG average has $\LambdaPG = 251 \pm 9~\mathrm{MeV}$ (${\sim}3.6\%$). A reduction to $\pm 2~\mathrm{MeV}$ by $\sim 2030$ would either confirm Match~1 at $< 1\%$ or pull the central value $> 5\sigma$ from $6\pi/5$, falsifying it. \emph{Prediction}: if $\LambdaPG$ converges outside $251 \pm 2~\mathrm{MeV}$ at $< 2\%$ precision, the pure-gauge fingerprint fails.

\textit{Route~2 --- lattice extension to non-$\SU(3)$ gauge groups.}\quad The framework gives $\kappa(\SO(5)) = 1/8$ (since $|\Phiroots(B_2)| = 4$), $\kappa(G_2) = 1/12$, $\kappa(\mathrm{Sp}(2)) = 1/8$. Hence $\Piskappa(\SO(5)) = 8\pi/7$, $\Piskappa(G_2) = 12\pi/11$. If lattice extensions to glueball-binding-ratio observables in $\SO(5)$ or $G_2$ gauge theory~\cite{AthenodorouTeper2021} match these $G$-specific values at $\le 2\%$, the framework gains substantial support; if they fail, the fingerprint is restricted to $\SU(3)$.

\textit{Route~3 --- $V_{ud}$ precision.}\quad Targets of $\pm 0.00003$ by $\sim 2030$ from electroweak radiative corrections~\cite{SengGorchtein2018} and super-allowed CKM unitarity tests will probe the $\sim 10\sigma$ discrepancy with $35/36$. A move toward $0.97222$ confirms; a move away (e.g.\ toward $0.9740$) falsifies.

\textit{Route~4 --- magnetic-moment precision.}\quad Belle~II / STCF programs target ${\sim}0.1\%$ precision on hyperon properties by 2030. A central $|\mu_{\Sigma^+}/\mu_{\Xi^-}|$ moving more than $\pm 0.04$ from $6\pi/5$ would falsify Match~3.

\textit{Combined timeline.}\quad Routes~1, 3, 4 are independently testable at $< 1\%$ by $\sim 2030$. Route~2 is testable at any time given dedicated lattice computation on non-$\SU(3)$ ensembles. A coordinated 5-year programme would either confirm the framework's $\kappa$-fingerprint at ${\sim}4\sigma$ aggregate or rule it out at $> 5\sigma$.

\section{Honest scope and disclaimers}
\label{sec:scope}

\textit{What this paper is.}\quad A \emph{numerical-pattern observation paper}, methodologically analogous to the Koide formula~\cite{Koide1983} (lepton-mass identity at $< 10^{-4}$ relative precision, unexplained for 43 years), the Titius--Bode relation (1772), or the Veneziano amplitude (1968). Such patterns sometimes precede genuine derivations and sometimes turn out to be accidents; our four-fold $7.75\sigma$ excess inside $\calF_\kappa$ is evidence, not proof.

\textit{What this paper is not.}\quad
\begin{itemize}
\item \emph{It is not a derivation.} None of the four matches is derived from first principles; we do not claim to predict $m_p$, $\langle\bar{q}q\rangle$, $\mu_B$, or $V_{ud}$.
\item \emph{It is not a unification.} The three sectors (QCD bound states, hadronic magnetic moments, CKM) are not related by any derivation in this work.
\item \emph{It is not free of human selection.} The formula library $\calF$ was constructed by the author with prior knowledge of $\kappa = 1/6$. The $17.09\sigma$ value is conditional on $\calF$; alternative libraries give different numerical values.
\end{itemize}

\textit{Logical independence.}\quad The structural companion papers~\cite{Remondiere2026a,Remondiere2026b} are logically independent of the present empirical observations. Falsification of the present pattern does not invalidate the structural results, which rest on Lie-algebraic and Lean-certified arguments. Conversely, full confirmation augments the structural results with an empirical bridge to phenomenology but does not constitute a derivation.

\section*{Acknowledgments and COPE-compliant disclosure}
\label{sec:ack}

The author thanks the FLAG~2024 working group for the consensus $\LambdaPG$ value and the PDG~2024 group for the magnetic-moment and CKM averages. All numerical values were extracted by direct reading of the original publications, with arXiv IDs verified against the live arXiv API to prevent citation fabrication.

Per COPE recommendations on AI-assisted research \cite{COPE2023}: AI tools (Anthropic Claude Opus~4.x and DeepSeek V4-Pro, 2026-05) were used as adversarial-review assistants for arXiv reference cross-checks, Python / NumPy Bonferroni computation, and scheme-dependence ambiguity identification. AI tools did not author and are not credited as authors; all intellectual content was generated and is owned by the human author. Numerical reproducibility: every value is reproducible from the Python scripts archived at \texttt{[Zenodo DOI to be assigned upon submission]}.

The author is grateful to anonymous referees in advance for adversarial criticism, particularly on the question of whether the Bonferroni baseline of \S~\ref{sec:bonferroni} adequately controls for dependence between rational-formula candidates.

\begin{thebibliography}{99}

\bibitem{PDG2024}
R.~L.~Workman et al.\ (Particle Data Group),
\emph{Review of Particle Physics},
Prog.\ Theor.\ Exp.\ Phys.\ \textbf{2024}, 083C01 (2024);
\href{https://pdg.lbl.gov}{pdg.lbl.gov}.

\bibitem{FLAG2024}
Y.~Aoki et al.\ (FLAG Working Group),
\emph{FLAG Review 2024},
Eur.\ Phys.\ J.\ C \textbf{84}, 712 (2024);
\href{https://arxiv.org/abs/2411.04268}{arXiv:2411.04268}.

\bibitem{SchaeferSommerVirotta2011}
S.~Schaefer, R.~Sommer, F.~Virotta (ALPHA Collaboration),
\emph{Critical slowing down and error analysis in lattice QCD simulations},
Nucl.\ Phys.\ B \textbf{845}, 93 (2011);
\href{https://arxiv.org/abs/1009.5228}{arXiv:1009.5228}.

\bibitem{LuscherSchaefer2011}
M.~L\"uscher, S.~Schaefer,
\emph{Lattice QCD without topology barriers},
JHEP \textbf{1107}, 036 (2011);
\href{https://arxiv.org/abs/1105.4749}{arXiv:1105.4749}.

\bibitem{Koide1983}
Y.~Koide,
\emph{New view of quark and lepton mass hierarchy},
Phys.\ Rev.\ D \textbf{28}, 252 (1983);
\emph{Fermion-boson two-body model of quarks and leptons and Cabibbo mixing},
Lett.\ Nuovo Cim.\ \textbf{34}, 201 (1982).

\bibitem{Remondiere2026a}
K.~R\'emondi\`ere,
\emph{A five-condition uniqueness theorem for the gauge sector of the Standard Model},
manuscript, 2026; arXiv submission in preparation.

\bibitem{Remondiere2026b}
K.~R\'emondi\`ere,
\emph{Synthesis v2: log-Sobolev framework for compact-Lie-group gauge theory},
manuscript, 2026; Zenodo DOI to be assigned.

\bibitem{Alexandrou1812}
C.~Alexandrou, S.~Bacchio, M.~Constantinou, J.~Finkenrath, K.~Hadjiyiannakou, K.~Jansen, G.~Koutsou, A.~Vaquero Aviles-Casco,
\emph{Proton and neutron electromagnetic form factors from lattice QCD};
\href{https://arxiv.org/abs/1812.10311}{arXiv:1812.10311}.

\bibitem{HardyTowner2020}
J.~C.~Hardy, I.~S.~Towner,
\emph{Superallowed $0^+ \to 0^+$ nuclear $\beta$ decays: 2020 critical survey, with implications for $V_{ud}$ and CKM unitarity},
Phys.\ Rev.\ C \textbf{102}, 045501 (2020). [Journal-only; no arXiv preprint at the time of writing.]

\bibitem{SengGorchtein2018}
C.-Y.~Seng, M.~Gorchtein, H.~H.~Patel, M.~J.~Ramsey-Musolf,
\emph{Reduced hadronic uncertainty in the determination of $V_{ud}$},
Phys.\ Rev.\ Lett.\ \textbf{121}, 241804 (2018);
\href{https://arxiv.org/abs/1807.10197}{arXiv:1807.10197}.

\bibitem{SengGalviz2021}
C.-Y.~Seng, D.~Galviz, W.~J.~Marciano, U.-G.~Mei\ss{}ner,
\emph{Update on $|V_{us}|$ and $|V_{us}/V_{ud}|$ from semileptonic kaon and pion decays},
\href{https://arxiv.org/abs/2107.14708}{arXiv:2107.14708}.

\bibitem{Luscher2010}
M.~L\"uscher,
\emph{Properties and uses of the Wilson flow in lattice QCD},
JHEP \textbf{1008}, 071 (2010);
\href{https://arxiv.org/abs/1006.4518}{arXiv:1006.4518}.

\bibitem{AthenodorouTeper2021}
M.~Athenodorou, M.~Teper,
\emph{The glueball spectrum of $\SU(N)$ gauge theories in $D = 3 + 1$ dimensions},
JHEP \textbf{2021}, 172 (2021);
\href{https://arxiv.org/abs/2106.00364}{arXiv:2106.00364}.

\bibitem{Sommer1994}
R.~Sommer,
\emph{A new way to set the energy scale in lattice gauge theories and its application to the static force and $\alpha_s$ in $\SU(2)$ Yang--Mills theory},
Nucl.\ Phys.\ B \textbf{411}, 839 (1994);
\href{https://arxiv.org/abs/hep-lat/9310022}{arXiv:hep-lat/9310022}.

\bibitem{NeccoSommer2002}
S.~Necco, R.~Sommer,
\emph{The $N_f = 0$ heavy quark potential from short to intermediate distances},
Nucl.\ Phys.\ B \textbf{622}, 328 (2002);
\href{https://arxiv.org/abs/hep-lat/0108008}{arXiv:hep-lat/0108008}.

\bibitem{COPE2023}
Committee on Publication Ethics,
\emph{Authorship and AI tools},
COPE position statement, 2023;
\href{https://publicationethics.org/cope-position-statements/ai-author}{publicationethics.org/cope-position-statements/ai-author}.

\end{thebibliography}

\end{document}
